%% kapa: exam=1 round=35 year=2024 period=1 form=A subject=real_estate subject_ko=부동산학원론 seq=29 number=109
%% source: Q-Net 감정평가사 기출문제 (cst00309) artlSeq=5213229 file=raw/1008/5213229/1교시 문제지 원본.pdf
%% conversion: kordoc 4.24.51 -> tools/gen1.py (auto); review=proofread-sonnet; flags=none
\begin{question}{109}
아파트시장에서 아파트의 수요곡선을 우측(우상향)으로 이동시킬 수 있는 요인은 모두 몇 개인가? (단, 다른 조건은 동일함)
\begin{framedtext}
\begin{tabular}{@{}p{0.45\linewidth}p{0.5\linewidth}@{}}
○ 아파트 가격의 하락 & ○ 대체 주택 가격의 상승 \\
○ 총부채원리금상환비율(DSR) 규제 완화 & ○ 가구수 증가 \\
○ 모기지 대출(mortgage loan) 금리의 상승 & ○ 수요자의 실질 소득 감소 \\
○ 부채감당률(DCR) 규제 강화 \\
\end{tabular}
\end{framedtext}
\begin{choices}
  \item 2개
  \item 3개
  \item 4개
  \item 5개
  \item 6개
\end{choices}
\end{question}
\answer{2}
